%
% Solver
% ------------
\section{Nozzle CFD}

Throughout this document the test case is referred to as a {\it nozzle}.
This is the more common terminology for transonic flows. 
For incompressible flows, {\it channel} is more commonly used, but we use a
single terminology for both flow regimes.

\begin{figure}[ht]
  \begin{center}
    \includegraphics[width=0.70\textwidth]{c128-sol-simple.png}
  \end{center}
  \caption{Solution for compressible nozzle with 128 stations.}
  \label{fig-c128-sol-simple}
\end{figure}

\Cref{fig-c128-sol-simple} shows the converged flow solution for a transonic 
compressible test case.
The bottom left figure shows the area of the channel, $A$, which is derived from the
isentropic flow relationships:

\begin{equation}
  \frac{A}{A^*} = \frac{1}{M}
                  \left[ \frac{2}{\gamma+1} \left(1 + \frac{\gamma - 1}{2}M^2 \right)
                  \right]^{\frac{\gamma+1}{2(\gamma - 1)}}
  \label{eqn:isent-area}
\end{equation}

where $M$ is the Mach number, $\gamma$ is the ratio of specific heats and
$A^*$ is the area at the throat where $M=1$.
In this test case, all variables are non-dimensionalised so that $A^*=1$.
The length of the nozzle is always set to 1 and the user specifies an inlet Mach number, $M_{in}$.
To set-up the nozzle area distribution, the outlet Mach number is set to $M_{out} = 1+M_{in}$.
The axial stations are set according an equal Mach number of increment of
$\delta M = 1.0/(nx-1)$ where $nx$ is the number of axial stations in the mesh. 
The area at each station is set from \Cref{eqn:isent-area}.
Note that the area equation is for isentropic, i.e. reversible, flow where the Mach number
increases linearly. In the absence of a shock wave this gives the CFD solution.
However, the interest is in the case when there is a shock wave within the nozzle as shown 
in the bottom right of \Cref{fig-c128-sol-simple}.

For incompressible flows, the channel area is an arc of a circle with a specified
area contraction, typically $10\%$ at the midpoint, relative to the inlet and exit areas
which are both set to $1.0$. The area distribution is not shown but follows the shape
of the pressure distribution.

\begin{figure}[ht]
  \begin{center}
    \includegraphics[width=0.70\textwidth]{i128-sol-simple.png}
  \end{center}
  \caption{Solution for incompressible nozzle with 128 stations.}
  \label{fig-i128-sol-simple}
\end{figure}


%
% Boundary conditions
% -------------------
\subsection{Boundary Conditions}

In addition to setting the geometry of the nozzle, we also need to provide boundary
conditions at the entrance and exit to the nozzle. These must be specified in a way that
is consistent with that signals propogate within the flow.

Following \Cref{fig-nozzle-disc}, a velocity boundary condition is set at the inlet as
the velocity updates are largely convected in the flow direction.
a pressure boundary condition is set at the outlet, as pressure signals propogate in 
all directions at the speed of sound, relative to the flow field. 
Setting an exit pressure boundary condition is only physically valid when the exit 
flow velocity is subsonic. This is case for both incompressible and compressible flows
used here and, hence, both have the same form of outlet condition:

\begin{equation}
  p|_{i=n}  = p_{exit}
  \label{eqn-inletbc}
\end{equation}

That is, $p$ at the exit station, $i=n$ has a fixed value, $p_{exit}$.
Although, incompressible and compressible flows have the same form of exit boundary
condition, the values are different.
The compressible solution in \Cref{fig-c128-sol-simple} has $p_{exit} = 0.75$.
Whereas, the incompressible solution in \Cref{fig-i128-sol-simple} 
has $p_{exit} = 0.95$.

For both compressible and incompressible flows, the values of velocity at the exit and
the pressure at the inlet are computed from the solution of the flow equations.
At the inlet, the computed inlet value of pressure is used to set an inlet velocity 
according to the inflow boundary condition. 

\subsubsection{Compressible inflow condition}

For compressible flows, the isentropic relationships give:

\begin{equation}
  \left( \frac{T_0}{t} \right ) =  \left( \frac{P_0}{p} \right ) ^{\frac{\gamma-1}{\gamma}}
  \label{eqn-t-p-isen-1}
\end{equation}

Where $T_0$ and $P_0$ are the staganation temperature and pressure which are constant in
isentropic, i.e. shock-free, flows.
The flow upstream of the nozzle inlet is taken as isentropic and $T_0$ and $P_0$ 
are user-specified boundary values. 
For non-dimensionalised conditions, as used here, it is usual to set them both to unity.

Rearranging \Cref{eqn-t-p-isen-1} gives:
\begin{equation}
  t = T_0 \left( \frac{p}{P_0} \right ) ^{\frac{\gamma-1}{\gamma}}
  \label{eqn-t-p-isen-2}
\end{equation}

In addition, we have conservation of energy which, for isentropic flow, can be 
expressed in terms of temperature as:

\begin{equation}
  T_0 = t + \frac{1}{2c_p}u^2
  \label{eqn-energy}
\end{equation}

And can be rearranged to give:

\begin{equation}
  u = \sqrt{2 c_p (T_0 - t)}
  \label{eqn-uin-isen}
\end{equation}

\Cref{eqn-t-p-isen-2} and \Cref{eqn-uin-isen} allow the inlet velocity to be updated from
the value of pressure and the specified inlet stagnation conditions.


\subsubsection{Incompressible inflow condition}

For incompressible flow, we cannot use the isentropic relations as strictly the Mach number is
zero or, more accurately, the speed of sound is infinite. Instead, we uswe Bernoulli's equation:

\begin{equation}
  P_0 = p + \frac{1}{2}\rho u^2
  \label{eqn-Bernoulli}
\end{equation}

From which, the inlet velocity can be set from the computed inlet static pressure and
the specified inlet stagnation pressure:

\begin{equation}
  u = \sqrt{\frac{2}{\rho} (P_0 - p)}
  \label{eqn-uin-bern}
\end{equation}



% Discretisation
% --------------
\subsection{CFD discretisation}
\label{subsec-discrete}

Since, the focus of this test case is not CFD development, we use the discretisation scheme of
\cite{gupta1993development}. In this scheme, the velocities and pressures are stored at the 
same set of nodes with upwinded momentum control volumes and downwinded continuity control
volumes. These are show in \Cref{fig-nozzle-v-disc} and \Cref{fig-nozzle-disc} respectively.

\begin{figure}[h]
  \centering
  \begin{subfigure}[b]{0.70\textwidth}
      \centering
      \includegraphics[width=\textwidth]{nozzle-v-disc.png}
      \caption{Upwinded momentum control volumes used to update nodes 2,..,n.}
      \label{fig-nozzle-v-disc}
  \end{subfigure}
  \hfill
  \begin{subfigure}[b]{0.70\textwidth}
      \centering
      \includegraphics[width=\textwidth]{nozzle-p-disc.png}
      \caption{Downwinded continuity control volumes used to update nodes 1,..,n-1.}
      \label{fig-nozzle-p-disc}
  \end{subfigure}
  \caption{Control volume setting for momentum and continuity equations.}
  \label{fig-nozzle-disc}
\end{figure}

The 1-dimensional steady state equations to be solved are:

\begin{equation}
  \begin{array}{rcl}
    \rho u \frac{\partial u}{\partial x} &=& - \frac{\partial p}{\partial x} \\[8pt]
    \frac{\partial \rho u}{\partial x} &=& 0 \\
  \end{array}
  \label{eqn-euler-1d}
\end{equation}

For compressible flows, the following equations are also applied:

\begin{equation}
  \begin{array}{rcl}
    p   &=& R\rho t \\[8pt]
    T_0 &=& t + \frac{1}{2c_p}u^2
  \end{array}
  \label{eqn-gaslaw}
\end{equation}

where $R$ is the gas constant, $c_p$ is the specific heat at constant pressure, and
$T_0$ is the constant total temperature which is set in $1.0$ in non-dimensional 
units. This equation represents the conservation of energy. 


The differential equations are discretised on the 1-dimensional mesh and 
updated iteratively via corrections:
\begin{equation}
  \begin{array}{rcl}
     \ket{u}    &\gets& \ket{u}    + \ket{\delta u} \\[8pt]
     \ket{p}    &\gets& \ket{p}    + \ket{\delta p} \\[8pt]
     \ket{\rho} &\gets& \ket{\rho} + \ket{\delta \rho}
  \end{array}
  \label{eqn-iter-update}
\end{equation}

For incompressible cases, the density is fixed at a non-dimensional value of $1.0$ 
and is not updated as part of the solution. 

The coupled iterative scheme is the most conceptually straightforward and will
be described first


% Incompressible Coupled
% ----------------------
\subsubsection{Incompressible coupled solver}
\label{subsubsec-inco-couple}

For incompressible flow, the discrete version of momentum equation is \Cref{eqn-euler-1d} is:

\begin{equation}
  \begin{array}{rcl}
    \rho u_{i-\frac{1}{2}} \frac{u_i - u_{i-1}}{x_i - x_{i-1}} &=&
                   - \frac{p_i - p_{i-1}}{x_i - x_{i-1}} \\
  \end{array}
  \label{eqn-euler-1d-disc1}
\end{equation}

Note that we have followed \cite{gupta1993development} in using a finite difference
discretisation for the momentum equation.


It is common practice in CFD to form an {\it inner} linearised system
to solve for the corrections. The only non-linear term in \Cref{eqn-euler-1d-disc1} is
the term on the left hand side of the first equation. 
Following \cite{gupta1993development},
this is linearised using a fixed point iterative scheme for the momentum equation to give:

\begin{equation}
  \begin{array}{rcl}
    \rho u_{i-\frac{1}{2}} (u_i + \delta u_i - u_{i-1} - \delta u_{i-1}) &=&
                   - (p_i + \delta p_i - p_{i-1} - \delta p_{i-1} ) \\[8pt]
    \rho u_{i-\frac{1}{2}} (\delta u_i - \delta u_{i-1}) +
    (\delta p_i - \delta p_{i-1}) &=&
    -\rho u_{i-\frac{1}{2}} (u_i - u_{i-1}) -
    (p_i - p_{i-1}) \\[8pt]
    &=& R_i^u
  \end{array}
  \label{eqn-euler-1d-disc2}
\end{equation}

We have also followed \cite{gupta1993development} in using a finite difference
discretisation for the momentum equation.

For the continuity equation, we use a finite volume discretisation to give
\begin{equation}
  \begin{array}{rcl}
    A_{i+1}\rho u_{i+1} - A_{i}\rho u_{i} &=& 0 \\[8pt]
    A_{i+1}\rho \delta u_{i+1} - A_{i}\rho \delta u_{i} &=&
    -(A_{i+1}\rho u_{i+1} - A_{i}\rho u_{i}) \\[8pt]
    &=& R_i^p
  \end{array}
  \label{eqn-euler-1d-disc3}
\end{equation}

Where $R_u^i$ and $R_p^i$ are the errors in satisfying the flow equations at node $i$,
these are usually called the residuals.

We must also linearise the incompressible inlet boundary condition from \Cref{eqn-Bernoulli}:

\begin{equation}
  \begin{array}{rcl}
    p_1 + \delta p_1 + \frac{1}{2} \rho (u_1 + \delta u_1) ^2 &=& P_0 \\[8pt]
          \delta p_1 + \rho u_1 \delta u _1&=& P_0 - p_1 - \frac{1}{2} \rho u_1^2 \\[8pt]
          \delta p_1 + \rho u_1 \delta u_1 &=& R^u_1
  \end{array}
  \label{eqn-incomp-in-disc1}
\end{equation}

The exit boundary condition is simply:
\begin{equation}
  \delta p_n = p_{exit} - p_n
  \label{eqn-incomp-ex-disc1}
\end{equation}


Collecting together the discrete momentum and continuity equations leads to a set of
simultaneous equations that can be expressed in block matrix form as:


\begin{equation}
  \begin{pmatrix}
    A^u + B^{u_1} & D^x + B^{p_1} \\
    M^x           & B^{p_n} \\
  \end{pmatrix}
  \begin{pmatrix}
    \ket{\delta u} \\
    \ket{\delta p}
  \end{pmatrix}
  = 
  \begin{pmatrix}
    \ket{R^u} \\
    \ket{R^p}
  \end{pmatrix}
  \label{eqn-implicit-01}
\end{equation}

Where
\begin{itemize}
    \item $A^u$ is the matrix containing the discrete convection operator 
                for the $u$ momentum equation.
    \item $B^{u_1}$ and $B^{p_1}$ are the the entries for the discretied inlet velocity boundary condition.
    \item $D^x$ is the discrete pressure gradient operator.
    \item $M^x$ is the the discrete continuity operator.
    \item $B^{p_n}$ is the the entry for the discretised outlet pressure boundary condition.
    \item $R^u$ and $R^p$ are the residuals of the momentum and continuity equations.
\end{itemize}

\Cref{fig-i16-cpl-sparse} shows the sparsity pattern for an implicit coupled incompressible
matrix with 16 mesh nodes.
Row 0 corresponds to the inlet boundary condition.
Rows 1-15 are the momentum equations up to the exit node.
Rows 16-30 are the continuity equations starting at the inlet node.
Row 31 is the exit boundary condition.

\begin{figure}[h]
  \begin{center}
    \includegraphics[width=0.60\textwidth]{i16_cpl_sparse.png}
  \end{center}
  \caption{Incompressible nozzle with 16 stations, sparsity pattern for coupled matrix.}
  \label{fig-i16-cpl-sparse}
\end{figure}


\Cref{tab-imp-incomp} provides details for a range of mesh sizes for
incompressible implicit matrices. The dimensions of the Hermitian matrix are do the symmetrised
matrix needed by HHL. The number of unitaries are derived from the decomposition of
the Hermitian matrix into Pauli strings.
The number of CFD iterations are those required to reduce the residuals and 
corrections to less than $1 \times 10^{-12}$.

\begin{table}[h]
  \centering
  \begin{tabular}{c c c c c c c}
    \toprule
    Nstations & Hermitian matrix &  $\lambda_{min}$ & $\lambda_{max}$ & $\kappa$ & Nunitaries & CFD iterations \\
    \midrule
     8   &   32x32     &$7.9\times10^{-2}$ & 3.32      & 42.1     & 128   &  5  \\
     16  &   64x64     &$3.9\times10^{-2}$ & 3.36      & 86.1     & 316   &  5  \\
     32  &  128x128    &$1.9\times10^{-2}$ & 3.38      & 173      & 740   &  5  \\
     64  &  256x256    &$9.6\times10^{-3}$ & 3.39      & 348      & 1,752 &  5  \\
    \bottomrule
  \end{tabular}
  \caption{Incompressible implicit matrices.
           Dimensions, eigenvalue range, condition number and number of unitaries in LCU.
           Sampled after at end of run. Area contraction = 10\%.}
  \label{tab-imp-incomp}
\end{table}



% Compressible Coupled
% --------------------
\subsubsection{Compressible coupled solver}
\label{subsubsec-comp-couple}

For compressible flows, the implicit discretisation combines the gas law and energy equations
to give:

\begin{equation}
  \begin{array}{rcl}
    p_i + \delta p_i &=& R (\rho_i + \delta \rho_i)(T_0 - \frac{1}{2c_p}(u_i + \delta u_i)^2) \\[8pt]
                     &=& R\rho_i(T_0 - \frac{1}{2c_p}(u_i^2 + 2 u_i \delta u_i)) +
                         R \delta \rho_i (T_0 - \frac{1}{2c_p}u_i^2) \\[8pt]
                     &=& R\rho_i t_i + \frac{R}{c_p} \rho_i u_i \delta u_i + Rt_i \delta \rho_i 
  \end{array}
  \label{eqn-euler-1d-disc4}
\end{equation}

Collecting the terms together gives:
\begin{equation}
  \begin{array}{rcl}
   -\frac{R}{c_p} \rho_i u_i \delta u_i + \delta p_i - Rt_i \delta \rho_i &=& R\rho_i t_i - \delta p_i \\[8pt]
      &=& R_{\rho}^i
  \end{array}
  \label{eqn-euler-1d-disc5}
\end{equation}

There is one further aspect to the implementation of the gas law which is due to the fact
that the continuity control volumes are {\it downwinded} as shown in \Cref{fig-nozzle-p-disc}.
To maintain stability at supersonic Mach numbers the perfect gas must compensate for this by
using an upwinded effective pressure:

\begin{equation}
    \rho = \frac{1}{Rt} p_{eff}
  \label{eqn-peff1}
\end{equation}

For Mach numbers less than 2, this is discretised as \cite{gupta1993development}:

\begin{equation}
   p_{eff,i} = p_{i-1} + a_0(p_i - p_{i-1}) + \frac{a_1}{2} (p_i - p_{i-2})
  \label{eqn-peff2}
\end{equation}

Where
\begin{equation}
  \begin{array}{rcl}
    a_0 &=& \frac{0.8}{3} (\frac{4}{M^2} - 1) \\[8pt]
    a_1 &=& 1.0 - a_0
  \end{array}
  \label{eqn-peff3}
\end{equation}

Where $M$ is the local maximum Mach number between stations $i$ and $i-1$.
For low subsonic flow, $a_0=1$ and $a_1 = 01$ giving $p_{eff,i} = p_i$, i.e. no upwinding.

For the implicit discretisation, only the first order upwinding correction is retained.
Hence, \Cref{eqn-euler-1d-disc5} becomes:

\begin{equation}
   -\frac{R}{c_p} \rho_i u_i \delta u_i + \delta p_i - Rt_i \delta \rho_i = 
    R\rho_i t_i - \delta p_{i-1} - a_0(\delta p_i - \delta p_{i-1}) \\
  \label{eqn-euler-1d-disc6}
\end{equation}

We must also linearise the compressible inlet boundary condition from \Cref{eqn-energy}:

\begin{equation}
  \begin{array}{rcl}
    t_1 + \delta t_1 + \frac{1}{2c_p} (u_1 + \delta u_1) ^2 &=& T_0 \\[8pt]
    t_1 + \delta t_1 + \frac{1}{c_p} u_1 \delta u_1 &=& T_0 - \frac{1}{2c_p} \rho u_1^2 \\[8pt]
    t_1 + \delta t_1 + \frac{1}{c_p} u_1 \delta u_1 &=& R^u_1
  \end{array}
  \label{eqn-comp-in-disc1}
\end{equation}

Linearising \Cref{eqn-t-p-isen-2} gives:
\begin{equation}
  \begin{array}{rcl}
    t_1 + \delta t_1 &=& T_0 \left( \frac{p_1}{P_0} \right ) ^{\frac{\gamma-1}{\gamma}}
                             \left(1 + \frac{\delta p_1}{p_1} \right ) ^{\frac{\gamma-1}{\gamma}} \\[8pt]
                     &=& T_0 \left( \frac{p_1}{P_0} \right ) ^{\frac{\gamma-1}{\gamma}}
                             \left(1 + \frac{\gamma-1}{\gamma}\frac{\delta p_1}{p_1} \right )
  \end{array}
  \label{eqn-t-p-isen-3}
\end{equation}

Combining \Cref{eqn-comp-in-disc1} and \Cref{eqn-t-p-isen-3} gives:
\begin{equation}
  \begin{array}{rcl}
     \frac{1}{c_p} u_1 \delta u_1 + 
     \frac{T_0}{p_1} \left( \frac{p_1}{P_0} \right ) ^{\frac{\gamma-1}{\gamma}} \frac{\gamma-1}{\gamma}\delta p_1  &=&
     T_0 -T_0 \left( \frac{p_1}{P_0} \right ) ^{\frac{\gamma-1}{\gamma}} - \frac{1}{2c_p} \rho u_1^2 \\[8pt]
     \frac{1}{c_p} u_1 \delta u_1 +
     \frac{T_0}{p_1} \left( \frac{p_1}{P_0} \right ) ^{\frac{\gamma-1}{\gamma}} \frac{\gamma-1}{\gamma}\delta p_1 &=& R^u_1
  \end{array}
  \label{eqn-comp-in-disc2}
\end{equation}

The exit boundary condition is the same as \Cref{eqn-incomp-ex-disc1}.


The fixed point iteration means that the compressible discretisation is the same as the 
compressible one except that the local density is used rather than the constant one.
The compressible continuity equation is straightforward application of corrections to the
velocity and density. This leads to a coupled compressible matrix of the form:

\begin{equation}
  \begin{pmatrix}
    A^u + B^{u_1} & D^x + B^{p_1} & 0        & 0 \\
    M^u           & B^{p_n}       & M^{\rho} & 0 \\
    L^u           & L^p           & L^{\rho} & 0 \\
    0             & 0             & 0        & I \\
  \end{pmatrix}
  \begin{pmatrix}
    \ket{\delta u}    \\
    \ket{\delta p}    \\
    \ket{\delta \rho} \\
    \ket{\psi}        \\
  \end{pmatrix}
  = 
  \begin{pmatrix}
    \ket{R^u} \\
    \ket{R^p} \\
    \ket{R^{\rho}} \\
    \ket{0} \\
  \end{pmatrix}
  \label{eqn-implicit-02}
\end{equation}

where the newly introduced terms are:
\begin{itemize}
    \item $L^u$, $L^p$ and $L^\rho$ are the coefficients from the gas law discretisation.
    \item $M^{\rho}$ are the additional terms from the dicretisation of $\rho$ in the 
                continuity equation.
    \item $I$ is a diagonal block that is required to extend the matrix so that it has a 
              full rank equal to a power of 2. $\ket{\psi}$ are the corresponding entries in
              the state vector and the matrix solution ensures $\ket{\psi} = \ket{0}$.
              In practice the diagonal entries in I are set to have the same order as the
              other diagonal entries so as not to skew the eigenvalues.
\end{itemize}

\Cref{fig-c16-cpl-sparse} shows the sparsity pattern for an implicit coupled compressible
matrix with 16 mesh nodes.

\begin{figure}[ht]
  \begin{center}
    \includegraphics[width=0.60\textwidth]{c16_cpl_sparse.png}
  \end{center}
  \caption{Compressible nozzle with 16 stations, sparsity pattern for coupled matrix.}
  \label{fig-c16-cpl-sparse}
\end{figure}

\Cref{tab-imp-comp} provides details for a range of mesh sizes for
compressible implicit  matrices. The dimensions of the Hermitian matrix are do the symmetrised
matrix needed by HHL. The number of unitaries are derived from the decomposition of
the Hermitian matrix into Pauli strings.
The number of CFD iterations are those required to reduce the residuals and 
corrections to less than $1 \times 10^{-12}$.

\begin{table}[h]
  \centering
  \begin{tabular}{c c c c c c c}
    \toprule
    Nstations & Hermitian matrix &  $\lambda_{min}$ & $\lambda_{max}$ & $\kappa$ & Nunitaries & CFD iterations \\
    \midrule
     8   &   64x64     &$1.9\times10^{-2}$ & 3.01      & 157      & 416   &  96    \\  
     16  &  128x128    &$8.9\times10^{-3}$ & 3.28      & 369      & 1,024 &  177   \\  
     32  &  256x256    &$3.7\times10^{-3}$ & 3.51      & 941      & 2,432 &  410   \\  
     64  &  512x512    &$1.4\times10^{-3}$ & 3.65      & 2,556    & 5,632 &  1,004 \\
    \bottomrule
  \end{tabular}
  \caption{Compressible implicit matrices. 
           Dimensions, eigenvalue range, condition number and unitaries in LCU.
           Sampled after 100 iterations. Pressure ratio = 0.75}
  \label{tab-imp-comp}
\end{table}

From the above derivation, it can be seen that there are additional highly non-linear 
terms in the compressible discretisation, particularly for the effective pressure.
Whilst it is possible to recursively linearise these, it tends to produce an unstable
discretisation. Indeed it is often necessary to apply relaxation to the initial coupled 
iterations and perform some pre-SIMPLE iterations to provide a better initial guess. 
The compressible implicit solver still provides a factor of 3-4 speed-up relative to the 
compressible SIMPLE solver, see \Cref{tab-imp-comp} and \Cref{tab-pc-comp},
but is much less effective than for incompressible flows.



% Incompressible SIMPLE
% ---------------------
\subsubsection{Incompressible SIMPLE solver}
\label{subsubsec-inco-simple}

We follow the discretisation implemented in Appendix D of \cite{gupta1993development}.
First of all we solve the momentum equation \Cref{eqn-euler-1d-disc2} using a Jacobi scheme:

\begin{equation}
  \delta u_i = \frac{R_u^i}{\rho u_{i-\frac{1}{2}}}
  \label{eqn-simple-inco-umom}
\end{equation}

All the $\delta u_i$ terms are evaluated at each station before being used to update
the velocities. This equation is solved on the classical device as the calculation
of $\delta u_i$ can be overlapped with the matrix assembly. Note that the pressures 
are not updated at this stage as they are done separately using the pressure correction 
equation. The SIMPLE algorithm falls into a broader class of {\it segregated} solves.

To solve the continuity equation, the momentum equation is assumed to have been solved,
allowing the velocity and pressure corrections to be related by:

\begin{equation}
  \begin{array}{rcl}
    \delta u_i &=& -\frac{1}{\rho u_{i-\frac{1}{2}}} (\delta p_i - \delta p_{i-1}) \\[8pt]
    \delta u_{i+1} &=& -\frac{1}{\rho u_{i+\frac{1}{2}}} (\delta p_{i+1} - \delta p_{i})
  \end{array}
  \label{eqn-simple-inco-dudp}
\end{equation}

Note that the first equation is based on the momentum equation for station $i$ and
the second equation on the one at for station $i+1$.
These equations can be substituted into the continuity equation \Cref{eqn-euler-1d-disc3}
at station $i$ to give

\begin{equation}
    -\frac{A_{i+1}}{u_{i+\frac{1}{2}}} (\delta p_{i+1} - \delta p_{i})
    +\frac{A_{i}}{u_{i-\frac{1}{2}}} (\delta p_{i} - \delta p_{i-1}) = R_p^i 
  \label{eqn-simple-inco-pceqn1}
\end{equation}

Setting $a_i = \frac{A_{i}}{u_{i-\frac{1}{2}}}$, \Cref{eqn-simple-inco-pceqn1} can
be written as:

\begin{equation}
    -a_{i}\delta p_{i-1} + (a_{i} + a_{i+1})\delta p_{i} - a_{i+1}\delta p_{i+1}
    = R_p^i 
  \label{eqn-simple-inco-pceqn2}
\end{equation}

This leads to a matrix equation of the form $A^p \ket{\delta p} = \ket{R^p}$
which has the sparsity pattern shown \Cref{fig-i16-sim-sparse}.


\begin{figure}[h]
  \begin{center}
    \includegraphics[width=0.60\textwidth]{i16_sim_sparse.png}
  \end{center}
  \caption{Incompressible nozzle with 16 stations, sparsity pattern for pressure correction matrix.}
  \label{fig-i16-sim-sparse}
\end{figure}

Once the pressure correction has been solved, $\ket{\delta p}$ is used to an update to the
velocity field using corrections according to \Cref{eqn-simple-inco-dudp}.
Since the continuity equation is linear for incompressible flows, this results in a 
velocity field that exact satisfies continuity. This step is also completed on the
classical device as it does not require a matrix solution. 

\Cref{tab-pc-incomp} provides details for a range of mesh sizes for 
incompressible pressure corrections matrices. 
The dimensions of the Hermitian matrix are do the symmetrised
matrix needed by HHL. The number of unitaries are derived from the decomposition of
the Hermitian matrix into Pauli strings.
The number of CFD iterations are those required to reduce the residuals and 
corrections to less than $1 \times 10^{-12}$.


\begin{table}[h]
  \centering
  \begin{tabular}{c c c c c c c}
    \toprule
    Nstations & Hermitian matrix &  $\lambda_{min}$ & $\lambda_{max}$ & $\kappa$ & Nunitaries & CFD iterations \\
    \midrule
     8   &   16x16     &$1.4\times10^{-1}$ & 3.18      & 23.3     & 24  & 194   \\
     16  &   32x32     &$3.5\times10^{-2}$ & 3.36      & 97.4     & 56  & 389   \\
     32  &   64x64     &$8.7\times10^{-3}$ & 3.56      & 410      & 128 & 765   \\
     64  &  128x128    &$2.1\times10^{-3}$ & 3.81      & 1,747    & 288 & 1,491 \\
    \bottomrule
  \end{tabular}
  \caption{Incompressible pressure correction matrices.
           dimensions, eigenvalue range, condition number and unitaries in LCU.
           Sampled after after 10 iterations. Area contraction = 10\%.}
  \label{tab-pc-incomp}
\end{table}



% Compressible SIMPLE
% -------------------
\subsubsection{Compressible SIMPLE solver}
\label{subsubsec-comp-simple}

The compressible pressure correction equation follows the derivation of the
incompressible equation except that the additional terms in $\delta \rho$ are 
replaced by $\delta p$ terms using the the gas law. 
We follow \cite{gupta1993development} and used a fully upwinded effective pressure
for the gas law.
This produces a pressure correction matrix with the same sparsity pattern as 
\Cref{fig-i16-sim-sparse}.


\Cref{tab-pc-comp} provides details for a range of mesh sizes for 
compressible pressure corrections matrices. 
The dimensions of the Hermitian matrix are do the symmetrised
matrix needed by HHL. The number of unitaries are derived from the decomposition of
the Hermitian matrix into Pauli strings.
The number of CFD iterations are those required to reduce the residuals and 
corrections to less than $1 \times 10^{-12}$.

\begin{table}[h]
  \centering
  \begin{tabular}{c c c c c c c}
    \toprule
    Nstations & Hermitian matrix &  $\lambda_{min}$ & $\lambda_{max}$ & $\kappa$ & Nunitaries & CFD iterations \\
    \midrule
     8   &   16x16     & $4.8\times10^{-1}$ & 6.89    & 14.3     & 32  & 424   \\
     16  &   32x32     & $1.9\times10^{-1}$ & 8.22    & 42.9     & 80  & 700   \\
     32  &   64x64     & $7.1\times10^{-2}$ & 9.25    & 131      & 192 & 1,449 \\
     64  &  128x128    & $3.0\times10^{-2}$ & 10.1    & 341      & 448 & 2,884 \\
    \bottomrule
  \end{tabular}
  \caption{Compressible pressure correction matrices.
           Dimensions, eigenvalue range, condition number and unitaries in LCU
           for the compressible nozzle.(sampled after 100 iterations).}
  \label{tab-pc-comp}
\end{table}

